\documentclass[aspectratio=169,10pt]{beamer}

\usepackage{iftex}
\providecommand{\dunenoirsetupfonts}{}
\newcommand{\DuneNoirLocalFontPath}{fonts}

\newcommand{\DuneNoirMonoFallback}{%
  \IfFontExistsTF{JetBrainsMono Nerd Font Propo}{
    \setmonofont{JetBrainsMono Nerd Font Propo}[Scale=0.92]
  }{
    \IfFontExistsTF{JetBrainsMono Nerd Font}{
      \setmonofont{JetBrainsMono Nerd Font}[Scale=0.92]
    }{
      \IfFontExistsTF{Liberation Mono}{
        \setmonofont{Liberation Mono}[Scale=0.92]
      }{
        \IfFontExistsTF{DejaVu Sans Mono}{
          \setmonofont{DejaVu Sans Mono}[Scale=0.92]
        }{
          \setmonofont{Courier New}[Scale=0.92]
        }
      }
    }
  }
}

\newcommand{\DuneNoirSansFallback}{%
  \IfFontExistsTF{IBM Plex Sans}{
    \setsansfont{IBM Plex Sans}
    \setmainfont{IBM Plex Sans}
  }{
    \IfFontExistsTF{Helvetica Neue}{
      \setsansfont{Helvetica Neue}
      \setmainfont{Helvetica Neue}
    }{
      \IfFontExistsTF{Liberation Sans}{
        \setsansfont{Liberation Sans}
        \setmainfont{Liberation Sans}
      }{
        \IfFontExistsTF{DejaVu Sans}{
          \setsansfont{DejaVu Sans}
          \setmainfont{DejaVu Sans}
        }{
          \setsansfont{Arial}
          \setmainfont{Arial}
        }
      }
    }
  }
}

\newcommand{\DuneNoirUseHelvetica}{%
  \renewcommand{\dunenoirsetupfonts}{%
    \IfFontExistsTF{Helvetica Neue}{
      \setsansfont{Helvetica Neue}
      \setmainfont{Helvetica Neue}
    }{
      \IfFontExistsTF{Helvetica}{
        \setsansfont{Helvetica}
        \setmainfont{Helvetica}
      }{
        \DuneNoirSansFallback
      }
    }
    \DuneNoirMonoFallback
  }%
}

\newcommand{\DuneNoirUseOpenSans}{%
  \renewcommand{\dunenoirsetupfonts}{%
    \IfFileExists{\DuneNoirLocalFontPath/OpenSans.ttf}{
      \setsansfont{Open Sans Custom}[
        Path=\DuneNoirLocalFontPath/,
        UprightFont={OpenSans.ttf},
        ItalicFont={OpenSans-Italic.ttf},
        BoldFont={OpenSans.ttf},
        BoldItalicFont={OpenSans-Italic.ttf},
        BoldFeatures={FakeBold=1.45},
        BoldItalicFeatures={FakeBold=1.45}
      ]
      \setmainfont{Open Sans Custom}[
        Path=\DuneNoirLocalFontPath/,
        UprightFont={OpenSans.ttf},
        ItalicFont={OpenSans-Italic.ttf},
        BoldFont={OpenSans.ttf},
        BoldItalicFont={OpenSans-Italic.ttf},
        BoldFeatures={FakeBold=1.45},
        BoldItalicFeatures={FakeBold=1.45}
      ]
    }{
      \DuneNoirSansFallback
    }
    \DuneNoirMonoFallback
  }%
}

\newcommand{\DuneNoirUseDosis}{%
  \renewcommand{\dunenoirsetupfonts}{%
    \IfFileExists{\DuneNoirLocalFontPath/Dosis.ttf}{
      \setsansfont{Dosis Custom}[
        Path=\DuneNoirLocalFontPath/,
        UprightFont={Dosis.ttf},
        ItalicFont={Dosis.ttf},
        BoldFont={Dosis.ttf},
        BoldItalicFont={Dosis.ttf},
        ItalicFeatures={FakeSlant=0.16},
        BoldFeatures={FakeBold=1.6},
        BoldItalicFeatures={FakeBold=1.6,FakeSlant=0.16}
      ]
      \setmainfont{Dosis Custom}[
        Path=\DuneNoirLocalFontPath/,
        UprightFont={Dosis.ttf},
        ItalicFont={Dosis.ttf},
        BoldFont={Dosis.ttf},
        BoldItalicFont={Dosis.ttf},
        ItalicFeatures={FakeSlant=0.16},
        BoldFeatures={FakeBold=1.6},
        BoldItalicFeatures={FakeBold=1.6,FakeSlant=0.16}
      ]
    }{
      \DuneNoirSansFallback
    }
    \DuneNoirMonoFallback
  }%
}

\newcommand{\DuneNoirUseKoHo}{%
  \renewcommand{\dunenoirsetupfonts}{%
    \IfFileExists{\DuneNoirLocalFontPath/KoHo-Regular.ttf}{
      \setsansfont{KoHo Custom}[
        Path=\DuneNoirLocalFontPath/,
        UprightFont={KoHo-Regular.ttf},
        BoldFont={KoHo-Bold.ttf},
        ItalicFont={KoHo-Italic.ttf},
        BoldItalicFont={KoHo-BoldItalic.ttf}
      ]
      \setmainfont{KoHo Custom}[
        Path=\DuneNoirLocalFontPath/,
        UprightFont={KoHo-Regular.ttf},
        BoldFont={KoHo-Bold.ttf},
        ItalicFont={KoHo-Italic.ttf},
        BoldItalicFont={KoHo-BoldItalic.ttf}
      ]
    }{
      \DuneNoirSansFallback
    }
    \DuneNoirMonoFallback
  }%
}

\DuneNoirUseHelvetica
\makeatletter
\input{../dune_dark_template_lab/beamerthemeDUNENoir.sty}
\makeatother

\ifPDFTeX
  \usepackage[T1]{fontenc}
  \usepackage[utf8]{inputenc}
\fi

\usepackage{amsmath,amssymb}
\usepackage{array,booktabs,tabularx}
\usepackage{colortbl}
\usepackage{graphicx}
\usepackage{hyperref}
\usepackage{tikz}
\usetikzlibrary{arrows.meta,calc,positioning,fit,shapes.geometric}

\title[DAPHNE-SC Development]{DAPHNE-SC Development}
\subtitle{Moving board control from Linux-only service logic to Linux + RPU supervision}
\author{Manuel Arroyave}
\date{Draft, June 2026}
\newcommand{\mydate}{Draft, June 2026}

\definecolor{okgreen}{HTML}{8FD6A9}
\definecolor{warnyellow}{HTML}{E7C547}
\definecolor{riskred}{HTML}{F28B82}
\definecolor{softblue}{HTML}{9BDCF2}
\definecolor{softviolet}{HTML}{BBA7F2}
\definecolor{softorange}{HTML}{F7A072}
\definecolor{linuxpanel}{HTML}{255A68}
\definecolor{rpupanel}{HTML}{2D6245}
\definecolor{dpspanel}{HTML}{7A4427}
\definecolor{warnpanel}{HTML}{6F5B15}

\setbeamercolor{panel rpu}{fg=noirtext,bg=rpupanel}
\setbeamercolor{panel linux}{fg=noirtext,bg=linuxpanel}
\setbeamercolor{panel dps}{fg=noirtext,bg=dpspanel}
\setbeamercolor{panel warn}{fg=noirtext,bg=warnpanel}
\setbeamercolor{panel risk}{fg=white,bg=riskred!80!black}
\AtBeginEnvironment{tabularx}{\color{noirtext}\arrayrulecolor{noirline}}
\AtBeginEnvironment{tabular}{\color{noirtext}\arrayrulecolor{noirline}}

\newcommand{\eyebrow}[1]{{\color{noircyan}\small\bfseries #1\par}}
\newcommand{\tinycite}[1]{{\color{noirmuted}\tiny Source basis: #1\par}}

\newcommand{\panelbox}[3][panel]{%
  \begin{beamercolorbox}[wd=\linewidth,sep=5pt,rounded=true]{#1}
    {\usebeamercolor[fg]{#1}\bfseries #2\par}
    \vspace{0.35mm}
    {\usebeamercolor[fg]{#1}\footnotesize #3\par}
  \end{beamercolorbox}
}

\newcommand{\fixedpanel}[4][panel]{%
  \begin{beamercolorbox}[wd=\linewidth,sep=5pt,rounded=true]{#1}
    \begin{minipage}[t][#2][t]{\dimexpr\linewidth-10pt\relax}
      \raggedright
      {\usebeamercolor[fg]{#1}\bfseries #3\par}
      \vspace{0.35mm}
      {\usebeamercolor[fg]{#1}\footnotesize #4\par}
    \end{minipage}
  \end{beamercolorbox}
}

\newcommand{\statcard}[3][panel alt]{%
  \begin{beamercolorbox}[wd=\linewidth,sep=5pt,rounded=true]{#1}
    {\usebeamercolor[fg]{#1}\bfseries\fontsize{14}{16}\selectfont #2\par}
    \vspace{0.35mm}
    {\usebeamercolor[fg]{#1}\scriptsize #3\par}
  \end{beamercolorbox}
}

\newcommand{\sectioncard}[2]{%
  \begin{frame}[plain]
    \begin{tikzpicture}[remember picture,overlay]
      \fill[noirpanelalt] (current page.south west) rectangle (current page.north east);
      \fill[duneorange] ($(current page.north west)+(0.9,-0.9)$) rectangle ++(1.8,0.08);
      \fill[noircyan] ($(current page.south east)+(-2.8,0.9)$) rectangle ++(1.8,0.08);
      \draw[noirline, line width=0.8pt] ($(current page.south west)+(0.9,1.1)$) -- ($(current page.south east)+(-0.9,1.1)$);
      \node[anchor=north west,text=noirmuted,font=\small\bfseries] at ($(current page.north west)+(0.9,-1.35)$) {SECTION};
      \node[anchor=north west,align=left,text=white,text width=0.9\paperwidth,
        font=\bfseries\fontsize{28}{31}\selectfont] at ($(current page.north west)+(0.8,-3.55)$) {#1};
      \node[anchor=north west,align=left,text=noirmuted,text width=0.91\paperwidth,
        font=\fontsize{13}{16}\selectfont] at ($(current page.north west)+(0.8,-5.2)$) {#2};
    \end{tikzpicture}
  \end{frame}
}

\tikzset{
  sysbox/.style={
    draw=noirline,
    fill=noirpanelalt,
    text=noirtext,
    rounded corners=4pt,
    align=center,
    text width=2.75cm,
    minimum height=0.92cm,
    inner sep=4pt,
    font=\scriptsize
  },
  syswide/.style={sysbox,text width=3.4cm},
  linuxbox/.style={syswide,draw=noircyan!85!white,fill=linuxpanel,text=noirtext},
  rpubox/.style={syswide,draw=noirmint,fill=rpupanel,text=noirtext},
  plbox/.style={syswide,draw=softviolet,fill=noirpanel,text=noirtext},
  dpsbox/.style={syswide,draw=softorange,fill=dpspanel,text=noirtext},
  daqbox/.style={syswide,draw=softblue,fill=noirpanel,text=noirtext},
  arrow/.style={-Latex,line width=0.82pt,draw=noircyan},
  warnarrow/.style={-Latex,line width=0.9pt,draw=duneorange},
  hardarrow/.style={-Latex,line width=0.95pt,draw=softorange},
  softline/.style={line width=0.7pt,draw=noirline,dashed},
  smallnote/.style={align=center,font=\tiny,text=noirmuted}
}

\newcommand{\stackDiagram}{%
\resizebox{0.98\linewidth}{!}{%
\begin{tikzpicture}[node distance=0.66cm and 0.55cm]
  \node[plbox] (pl) {PL / gateware\\registers and IO};
  \node[rpubox,right=of pl] (rpu) {RPU supervisor\\local state machine};
  \node[linuxbox,right=of rpu] (linux) {Linux board service\\protobuf endpoint};
  \node[linuxbox,right=of linux] (bridge) {OPC-UA bridge\\Ignition tags};
  \node[daqbox,below=0.9cm of linux] (daq) {DAQ client\\run configure};
  \node[dpsbox,below=0.9cm of rpu] (dps) {DPS / supply inhibit\\hard protection};
  \draw[arrow] (pl) -- (rpu);
  \draw[arrow] (rpu) -- (linux);
  \draw[arrow] (linux) -- (bridge);
  \draw[arrow] (daq) -- (linux);
  \draw[hardarrow] (dps) -- (rpu);
  \draw[hardarrow] (dps) -| (pl);
  \node[smallnote,below=0.18cm of rpu] {continues when Linux or network is unavailable};
  \node[smallnote,below=0.18cm of bridge] {operator views, alarms, history};
\end{tikzpicture}
}%
}

\newcommand{\networkDownDiagram}{%
\resizebox{0.95\linewidth}{!}{%
\begin{tikzpicture}[node distance=0.68cm and 0.64cm]
  \node[linuxbox] (ign) {Ignition / OPC-UA\\may be unreachable};
  \node[linuxbox,right=of ign] (linux) {Linux service\\may restart};
  \node[rpubox,right=of linux] (rpu) {RPU\\keeps permits local};
  \node[plbox,right=of rpu] (loads) {DAPHNE rails, bias,\\AFE enables};
  \node[dpsbox,below=0.9cm of rpu] (hard) {PLC / supply inhibit\\if required};
  \draw[softline] (ign) -- node[above,font=\tiny,text=noirmuted]{network loss} (linux);
  \draw[arrow] (linux) -- (rpu);
  \draw[hardarrow] (rpu) -- (loads);
  \draw[hardarrow] (hard) -- (rpu);
  \draw[hardarrow] (hard) -| (loads);
  \node[smallnote,below=0.18cm of loads] {no new remote authority is assumed};
\end{tikzpicture}
}%
}

\begin{document}

\begin{frame}
  \titlepage
\end{frame}

\begin{frame}[shrink=8]{The Sceptical Position}
\begin{columns}[T,onlytextwidth]
  \column{0.55\linewidth}
    \eyebrow{SHORT ANSWER}
    {\color{white}\bfseries\fontsize{16}{19}\selectfont
    The RPU is a strong option only for local, measurable, fast protection.
    It is not a replacement for detector-wide knowledge.\par}
    \vspace{3mm}
    \begin{itemize}
      \item Good fit: bias current, rail current, board temperature, AFE power state.
      \item Bad fit: amount of light in the cryostat, who enabled a calibration system, global detector state.
      \item External permits can be inputs, but the RPU should not invent them.
    \end{itemize}
  \column{0.41\linewidth}
    \panelbox[panel rpu]{Use the RPU for}{Fast local shutdown when local electrical or thermal evidence says the board is unhealthy.}
    \vspace{2mm}
    \panelbox[panel warn]{Do not claim}{That the RPU understands optical exposure or global detector safety by itself.}
\end{columns}
\end{frame}

\sectioncard{Main Question}{Which Linux failure modes can the RPU catch locally, and which ones remain outside the board?}

\begin{frame}[shrink=8]{Main Slide: Failure Modes Covered By The RPU}
\begin{columns}[T,onlytextwidth]
  \column{0.48\linewidth}
    \panelbox[panel rpu]{Linux hangs or restarts}{RPU sees stale Linux heartbeat but still samples bias current, rail current, and temperature. It can trip both bias enables or AFE power on a local fault.}
    \vspace{2mm}
    \panelbox[panel rpu]{Linux polling is stale}{RPU uses fresh local samples, sample-age checks, debounce, and hysteresis. It latches a fault if samples are stale or unhealthy.}
  \column{0.48\linewidth}
    \panelbox[panel rpu]{Network / OPC-UA is down}{Operators lose visibility, but the RPU keeps the local protection loop alive and refuses new remote authority.}
    \vspace{2mm}
    \panelbox[panel rpu]{Bad enable after fault}{RPU keeps the fault latched, rejects re-enable, and requires explicit clear plus controlled recovery.}
\end{columns}
\vspace{2mm}
\panelbox[panel warn]{Sceptical boundary}{This only works if the RPU has an independent monitor path and a real actuator path to remove bias or AFE power. Otherwise it is just another alarm.}
\end{frame}

\begin{frame}[shrink=8]{Concrete Local Trip Cases}
\begin{columns}[T,onlytextwidth]
  \column{0.32\linewidth}
    \statcard[panel rpu]{Bias current}{Switch off both SiPM biases when current leaves the healthy window and persists through debounce.}
  \column{0.32\linewidth}
    \statcard[panel rpu]{Temperature}{Remove bias and/or AFE power when board, regulator, or AFE temperature crosses the trip threshold.}
  \column{0.32\linewidth}
    \statcard[panel rpu]{Rail current}{Disable AFE or bias rail when current/power indicates an electrical fault.}
\end{columns}
\vspace{2mm}
\panelbox[panel alt]{Why RPU instead of Linux}{The RPU can keep running a small finite-state machine while Linux restarts, the network is down, the OPC-UA bridge is stale, or DAQ is not responding.}
\vspace{2mm}
\panelbox[panel warn]{Sceptical limit}{This is not a physics trigger or a light monitor. It reacts to local electrical and thermal consequences, not to the original global cause.}
\end{frame}

\begin{frame}[shrink=8]{What The RPU Cannot Know Locally}
\begin{columns}[T,onlytextwidth]
  \column{0.48\linewidth}
    \begin{block}{Not local information}
      \begin{itemize}
        \item light level inside the cryostat;
        \item laser, LCM, or flasher state unless wired in;
        \item cryogenics, rack, or detector-wide permit state unless wired in;
        \item whether high SiPM current is physics light, noise, damage, or a control mistake.
      \end{itemize}
    \end{block}
  \column{0.48\linewidth}
    \begin{block}{Local information}
      \begin{itemize}
        \item bias voltage and bias current;
        \item local rail current and voltage;
        \item local temperature;
        \item watchdogs, stale samples, and fault latches;
        \item external permit bits if they are provided as inputs.
      \end{itemize}
    \end{block}
\end{columns}
\vspace{1mm}
\panelbox[panel alt]{Therefore}{The honest claim is narrower: the RPU protects against locally observable electrical and thermal unhealthy states.}
\end{frame}

\begin{frame}{Example: Bias Current Guard}
\centering
\resizebox{0.95\linewidth}{!}{%
\begin{tikzpicture}[node distance=0.64cm and 0.58cm]
  \node[plbox] (mon) {Bias monitor\\I, V, timestamp};
  \node[rpubox,right=of mon] (filter) {RPU filter\\debounce + hysteresis};
  \node[rpubox,right=of filter] (state) {Latched fault\\no auto-clear};
  \node[dpsbox,right=of state] (off) {Switch off\\both bias enables};
  \node[linuxbox,below=0.9cm of state] (report) {Linux / Ignition\\alarm + history};
  \draw[arrow] (mon) -- (filter);
  \draw[arrow] (filter) -- (state);
  \draw[hardarrow] (state) -- (off);
  \draw[arrow] (state) -- (report);
  \node[smallnote,below=0.18cm of filter] {reject single-sample spikes};
  \node[smallnote,below=0.18cm of off] {conservative safe action};
\end{tikzpicture}}
\vspace{4mm}
\begin{columns}[T,onlytextwidth]
  \column{0.48\linewidth}
    \panelbox[panel rpu]{Trip condition}{Current outside the healthy SiPM range for long enough to survive debounce and threshold hysteresis.}
  \column{0.48\linewidth}
    \panelbox[panel warn]{Hard question}{Do we trust the local current monitor enough to remove bias without Linux confirmation?}
\end{columns}
\end{frame}

\begin{frame}[shrink=8]{Should It Switch Off Both Biases?}
\begin{columns}[T,onlytextwidth]
  \column{0.48\linewidth}
    \begin{block}{Argument for both}
      \begin{itemize}
        \item common-cause fault is possible;
        \item mapping may be ambiguous during commissioning;
        \item detector protection beats availability;
        \item a single bad bias path may indicate a shared supply or board problem.
      \end{itemize}
    \end{block}
  \column{0.48\linewidth}
    \begin{block}{Argument against both}
      \begin{itemize}
        \item unnecessary detector dead region if the fault is isolated;
        \item false trip rate matters during operations;
        \item one threshold may not fit all temperature and SiPM conditions;
        \item recovery becomes slower and more disruptive.
      \end{itemize}
    \end{block}
\end{columns}
\vspace{2mm}
\panelbox[panel alt]{Conservative proposal}{Start with both-bias shutdown for high-confidence, persistent, class-A faults. Keep per-bias shutdown as a later refinement only if monitors and mapping are proven.}
\end{frame}

\begin{frame}[shrink=8]{AFE Shutdown Is A Different Decision}
\begin{columns}[T,onlytextwidth]
  \column{0.48\linewidth}
    \panelbox[panel rpu]{Good RPU case}{AFE rail current, AFE power, or local temperature is high enough that continued operation could damage hardware.}
    \vspace{2mm}
    \panelbox[panel rpu]{Local action}{Disable AFE power or place the frontend in a defined safe state, then latch the reason.}
  \column{0.48\linewidth}
    \panelbox[panel warn]{Do not overreach}{Do not turn off AFEs because the RPU suspects too much light. It cannot know that locally unless the light system permit is wired in.}
    \vspace{2mm}
    \panelbox[panel dps]{Escalation}{If the required action must survive board firmware or RPU failure, it belongs in DPS or a supply-side inhibit.}
\end{columns}
\end{frame}

\sectioncard{The Algorithm Must Be Boring}{For protection, boring is a feature: thresholds, timers, latches, explicit reset, and testable failure behavior.}

\begin{frame}{What The Local Algorithm Should Be}
\begin{columns}[T,onlytextwidth]
  \column{0.48\linewidth}
    \begin{block}{Allowed}
      \begin{itemize}
        \item calibrated thresholds;
        \item hysteresis between trip and release;
        \item debounce time for noisy ADC samples;
        \item stale-sample detection;
        \item latched trip with explicit reset;
        \item no automatic re-enable of bias.
      \end{itemize}
    \end{block}
  \column{0.48\linewidth}
    \begin{block}{Not allowed}
      \begin{itemize}
        \item complex diagnosis of light source;
        \item ML or adaptive policy in the protection loop;
        \item hidden threshold changes;
        \item recovery without operator or DAQ revalidation;
        \item silent failure to a permissive state.
      \end{itemize}
    \end{block}
\end{columns}
\end{frame}

\begin{frame}[shrink=8]{What Must Be Proven Before We Trust It}
\small
\begin{tabularx}{\linewidth}{@{}>{\bfseries\color{noircyan}}p{0.27\linewidth}X@{}}
\toprule
Question & Required evidence \\
\midrule
Can the RPU really switch off bias? & Hardware path from RPU to bias enable or supply inhibit, tested under load. \\
\addlinespace
Are current and temperature measurements trustworthy? & Calibration, sample rate, stale-sample detection, and fault-injection tests. \\
\addlinespace
Are thresholds physically meaningful? & Healthy SiPM current ranges versus temperature, bias setting, noise, and expected operating state. \\
\addlinespace
What is the false-trip cost? & Operations impact if bias is removed during a good run. \\
\addlinespace
What is the missed-trip cost? & Damage or detector-risk analysis if the RPU waits too long or does nothing. \\
\addlinespace
What happens when the RPU fails? & Watchdog, fail-safe output state, Linux alarm, and DPS fallback where required. \\
\bottomrule
\end{tabularx}
\end{frame}

\begin{frame}[shrink=8]{Recommended Decision Language}
\begin{columns}[T,onlytextwidth]
  \column{0.52\linewidth}
    \eyebrow{SAY THIS}
    {\color{white}\bfseries\fontsize{15}{18}\selectfont
    We propose the RPU as a local protection layer for DAPHNE electrical and
    thermal faults, especially SiPM bias-current trips and AFE/rail thermal or
    over-current faults.\par}
    \vspace{3mm}
    \begin{itemize}
      \item It reacts locally when Linux or the network is unavailable.
      \item It latches faults and requires controlled recovery.
      \item It reports to Linux, Ignition, and DAQ after the fact.
    \end{itemize}
  \column{0.42\linewidth}
    \panelbox[panel warn]{Do not say}{The RPU protects against unknown light in the cryostat unless a light-system permit is wired to it.}
    \vspace{2mm}
    \panelbox[panel dps]{Do say}{Detector-wide or hard protection still belongs to DPS or supply-side inhibit when independence is required.}
\end{columns}
\end{frame}

\begin{frame}{Decision Gate}
\begin{columns}[T,onlytextwidth]
  \column{0.48\linewidth}
    \panelbox[panel rpu]{Proceed if}{Local bias current, rail current, and temperature monitors are reliable enough to justify automatic local shutdown.}
    \vspace{2mm}
    \panelbox[panel rpu]{Proceed if}{The RPU output path can remove both bias enables and, for thermal/electrical faults, place AFEs in a safe state.}
  \column{0.48\linewidth}
    \panelbox[panel warn]{Do not proceed if}{Thresholds are guesses, monitor freshness is not guaranteed, or recovery behavior is not defined.}
    \vspace{2mm}
    \panelbox[panel warn]{Do not proceed if}{The action must be independent of board software but is not wired through DPS or a supply-side inhibit.}
\end{columns}
\end{frame}

\begin{frame}[shrink=8]{Takeaway}
\begin{columns}[T,onlytextwidth]
  \column{0.55\linewidth}
    \eyebrow{BOTTOM LINE}
    {\color{white}\bfseries\fontsize{17}{20}\selectfont
    Yes, it makes sense, but only with a narrow claim:
    the RPU is a critical local guard for electrical and thermal health,
    not a detector-wide protection oracle.\par}
    \vspace{4mm}
    {\color{noirmuted}\small Bias current trip, temperature trip, rail-current trip, latch, report, controlled reset.\par}
  \column{0.41\linewidth}
    \panelbox[panel rpu]{Best RPU case}{Switch off both SiPM biases when local current is persistently outside the healthy range.}
    \vspace{2mm}
    \panelbox[panel dps]{Still needed}{DPS or supply inhibit for protection that must survive board-software failure.}
\end{columns}
\end{frame}

\end{document}
